% Universal semantic tensor-network TikZ grammar for TNLean.
%
% This file fixes the meaning of every glyph, wire, and region used throughout
% the repository.  A consumer may replace only the `tn theme ...` styles; the
% semantic `tn ...` styles and the atomic commands remain common to every
% document.
%
% Every index contraction is solid.  Virtual and physical contractions have
% different widths.  Morphism arrows, comparison arrows, annotations, and
% region boundaries are separate semantic marks and do not represent indices.
% Dashed curves indicate grouping, a basis change, or an identified periodic
% boundary.  Tensor labels do not change the glyph used for the tensor.
\usetikzlibrary{backgrounds,calc,decorations.pathreplacing,fit,matrix,positioning}

% The calculus uses a restrained semantic palette.  A theme changes these
% colour tokens, but never geometry, typography, or the meaning of a stroke.
\colorlet{tnInk}{black!92}
\colorlet{tnSecondaryInk}{black!55}
\colorlet{tnAccent}{red!65!black}
\colorlet{tnRegion}{blue!65!black}
\colorlet{tnSector}{violet!65!black}
\colorlet{tnPaper}{white}
\colorlet{tnTensorText}{white}
\colorlet{tnInsertionFill}{tnAccent!18!tnPaper}
\colorlet{tnInsertionText}{tnInk}

\tikzset{
    % Theme slots.  These are the only styles a light or dark document should
    % replace.  Semantic styles below never change meaning between themes.
    tn theme picture/.style={
            baseline=-0.1cm
        },
    tn theme tensor/.style={
            draw=tnInk,
            fill=tnInk,
            text=tnTensorText,
            shape=circle,
            inner sep=0.055cm
        },
    tn theme insertion/.style={
            draw=tnAccent,
            fill=tnInsertionFill,
            text=tnInsertionText,
            shape=circle,
            inner sep=0.055cm
        },
    tn theme junction/.style={
            fill=tnInk,
            shape=circle,
            inner sep=0.020cm
        },
    tn theme component/.style={
            draw=tnSecondaryInk,
            fill=tnPaper,
            inner sep=1.4pt
        },
    tn theme factor/.style={
            draw=tnSecondaryInk,
            fill=tnSecondaryInk!5!tnPaper,
            inner sep=1.4pt
        },
    tn theme map/.style={
            draw=tnSecondaryInk,
            rounded corners=2pt,
            fill=tnPaper,
            inner sep=2pt
        },
    tn theme state/.style={
            draw=tnSecondaryInk,
            rounded corners=2pt,
            fill=tnSecondaryInk!4!tnPaper,
            inner sep=2pt
        },
    tn theme expression/.style={
            draw=tnSecondaryInk,
            rounded corners=0.7pt,
            fill=tnPaper,
            inner sep=2pt
        },
    tn theme virtual wire/.style={draw=tnInk, line width=0.55pt},
    tn theme physical wire/.style={draw=tnInk, line width=0.55pt},
    tn theme trace/.style={rounded corners=3pt},
    tn theme factor region/.style={draw=tnSecondaryInk, rounded corners=2pt,
            line width=0.80pt},
    tn theme grouping region/.style={
            draw=tnSecondaryInk,
            densely dashed,
            rounded corners=2pt,
            line width=0.80pt
        },
    tn theme annotation/.style={draw=tnSecondaryInk, line width=0.40pt},
    tn theme label/.style={text=tnInk},
    % Semantic node classes.  Their names and meanings are theme-independent.
    tn picture/.style={tn theme picture},
    tn tensor node/.style={tn theme tensor},
    tn insertion node/.style={tn theme insertion},
    % A junction is a contraction point of index lines, not a local tensor.
    tn junction node/.style={tn theme junction},
    tn component node/.style={tn theme component},
    tn factor node/.style={tn theme factor},
    tn map node/.style={tn theme map},
    tn state node/.style={tn theme state},
    tn expression node/.style={tn theme expression},
    % Semantic index lines.
    tn virtual wire/.style={tn theme virtual wire},
    tn physical wire/.style={tn theme physical wire},
    tn virtual trace/.style={tn virtual wire, tn theme trace},
    tn physical trace/.style={tn physical wire, tn theme trace},
    tn bent wire/.style={tn virtual wire, tn theme trace},
    % Regions and non-contraction marks.  A factor boundary is not an index.
    tn factor region/.style={tn theme factor region},
    tn grouping region/.style={tn theme grouping region},
    tn annotation/.style={tn theme annotation},
    tn label/.style={tn theme label},
    tn basis change mark/.style={densely dashed, line width=0.40pt},
    tn periodic mark/.style={tn virtual wire, densely dashed, ->},
    tn morphism arrow/.style={draw=tnInk, ->, line width=0.55pt},
    tn comparison arrow/.style={draw=tnInk, <->, line width=0.55pt},
    tn selected path/.style={draw=tnAccent, line width=0.80pt},
    tn secondary path/.style={draw=tnRegion, line width=0.80pt},
    tn region selected/.style={
            draw=tnAccent,
            fill=tnAccent!12!tnPaper,
            line width=0.80pt,
            rounded corners=1mm
        },
    tn region secondary/.style={
            draw=tnRegion,
            fill=tnRegion!10!tnPaper,
            line width=0.80pt,
            rounded corners=1mm
        },
    tn region complement/.style={
            draw=tnSecondaryInk,
            fill=tnSecondaryInk!8!tnPaper,
            densely dashed,
            line width=0.80pt,
            rounded corners=1mm
        },
    tn region collar/.style={
            draw=tnSector,
            fill=tnSector!10!tnPaper,
            line width=0.80pt,
            rounded corners=1mm
        },
    tn edge distinguished/.style={line width=0.80pt, draw=tnAccent},
    tn vertex distinguished/.style={
            circle,
            draw=tnInk,
            fill=tnPaper,
            line width=0.80pt
        },
    tn region label selected/.style={text=tnAccent},
    tn region label secondary/.style={text=tnRegion},
    tn region label complement/.style={text=tnSecondaryInk},
    tn region label collar/.style={text=tnSector}
}

% Every semantic glyph is an instance of this one TikZ pic.  The public atom
% constructors choose the semantic style and declare the corresponding ports;
% the pic owns the canonical glyph geometry.
\tikzset{
    tn pic name/.store in=\TNpicname,
    tn pic style/.store in=\TNpicstyle,
    tn pic contents/.store in=\TNpiccontents,
    tn current pic options/.style={},
    pics/tn atom/.style={
        code={
            \node[\TNpicstyle,tn current pic options]
                (\TNpicname) at (0,0) {\TNpiccontents};
        }
    }
}

\makeatletter

% ---------- Semantic registry and event stream ----------

% Ports are stored in one picture-local property map.  The textual event
% stream is deliberately simple: it is both human-readable and stable under
% TeX-engine changes, and the repository audit canonicalizes it as a labelled
% graph after rendering.
\ExplSyntaxOn
\prop_new:N \g__tn_port_type_prop
\prop_new:N \g__tn_alias_source_prop
\iow_new:N \g__tn_event_iow

\AtBeginDocument
  { \iow_open:Nn \g__tn_event_iow { \c_sys_jobname_str.tnlog } }
\AtEndDocument { \iow_close:N \g__tn_event_iow }

\cs_new_protected:Npn \tn_event:n #1
  { \iow_now:Nx \g__tn_event_iow { #1 } }

\cs_new:Npn \tn_port_key:n #1
  { \int_use:N \TN@pictureid / \tl_to_str:e {#1} }

\cs_new_protected:Npn \tn_port_register:nn #1#2
  {
    \prop_if_in:NeTF \g__tn_port_type_prop { \tn_port_key:n {#1} }
      {
        \PackageError{tnlean-tikz}{Duplicate tensor-network port `#1'}
          {Every atom and alias must expose a distinct port name.}
      }
      {
        \prop_gput:Nee \g__tn_port_type_prop
          { \tn_port_key:n {#1} } {#2}
        \tn_event:n
          { port|picture=\int_use:N\TN@pictureid|name=\tl_to_str:e{#1}|type=#2 }
      }
  }

\cs_new_protected:Npn \tn_alias_register:nnn #1#2#3
  {
    \prop_if_in:NeTF \g__tn_alias_source_prop { \tn_port_key:n {#2} }
      {
        \prop_get:NeN \g__tn_alias_source_prop { \tn_port_key:n {#2} }
          \l_tmpa_tl
        \str_if_eq:eeT { \tl_to_str:e {#1} } { \tl_to_str:V \l_tmpa_tl }
          {
            \PackageError{tnlean-tikz}{Cyclic tensor-network alias `#1'}
              {Aliases must terminate at an atom port.}
          }
      }
      { }
    \tn_port_register:nn {#1} {#3}
    \prop_gput:Nee \g__tn_alias_source_prop
      { \tn_port_key:n {#1} } { \tl_to_str:e {#2} }
    \tn_event:n
      { alias|picture=\int_use:N\TN@pictureid|name=\tl_to_str:e{#1}|source=\tl_to_str:e{#2}|type=#3 }
  }

\cs_new_eq:NN \TNEventInternal \tn_event:n
\cs_new_eq:NN \TNPortRegisterInternal \tn_port_register:nn
\cs_new_eq:NN \TNAliasRegisterInternal \tn_alias_register:nnn

\cs_new_protected:Npn \tn_place_object:nn #1#2
  {
    \tl_if_eq:nnTF {#2} {at=origin}
      { \coordinate (#1Placement) at (0,0); }
      {
        \tl_if_in:nnTF {#2} {=}
          { \node[coordinate,#2] (#1Placement) {}; }
          { \coordinate (#1Placement) at #2; }
      }
  }
\cs_new_eq:NN \TNPlaceObjectInternal \tn_place_object:nn
\cs_new_protected:Npn \tn_placement_dispatch:nnn #1#2#3
  { \tl_if_in:nnTF {#1} {=} {#2} {#3} }
\cs_new_eq:NN \TNPlacementDispatch \tn_placement_dispatch:nnn
\ExplSyntaxOff

% Port declarations are local to one picture.  The registry itself is global
% so declarations made inside a TikZ loop remain visible after that loop ends.
\newcount\TN@pictureid
\tikzset{tn picture/.append style={
        execute at begin picture={\global\advance\TN@pictureid by 1\relax}
    }}

% Universal layout parameters.  Complete diagrams select a named layout
% profile; constructions use these names rather than repeat numerical
% distances.  A layer pitch is the displacement of one layer from the
% midline; paired layers therefore lie two pitches apart.
% A layout profile fixes all repeated local distances.  The normal profile is
% used for structural figures; compact is used for definitions and local
% identities.  Client diagrams select a profile by name and do not rescale or
% shift individual atoms.
\newcommand{\TN@normalprofile}{%
    \def\TN@layerpitch{0.50}%
    \def\TN@stackedseparation{1.00cm}%
    \def\TN@virtualleglen{0.60}%
    \def\TN@physicalleglen{0.42}%
    \def\TN@traceclearance{0.42}%
    \def\TN@busoffset{1.10}%
    \def\TN@buspitch{0.40}%
    \def\TN@branchfirstfraction{0.28}%
    \def\TN@branchsecondfraction{0.72}%
    \def\TN@cellpitch{1.10}%
    \def\TN@rowgap{0.52em}%
    \def\TN@relationgap{0.72em}%
    \def\TN@physlabeloff{0.68}%
    \def\TN@labelclearance{0.24}%
    \def\TN@annotationclearance{0.34}%
    \def\TN@chainsep{1.45}%
    \def\TN@chainright{4.20}%
    \def\TN@chainellipsisx{2.82}%
    \def\TN@gaugeradius{0.74}%
    \def\TN@gaugediagonal{0.54}%
}

\newcommand{\TN@compactprofile}{%
    \def\TN@layerpitch{0.40}%
    \def\TN@stackedseparation{0.80cm}%
    \def\TN@virtualleglen{0.48}%
    \def\TN@physicalleglen{0.34}%
    \def\TN@traceclearance{0.34}%
    \def\TN@busoffset{0.88}%
    \def\TN@buspitch{0.32}%
    \def\TN@branchfirstfraction{0.30}%
    \def\TN@branchsecondfraction{0.70}%
    \def\TN@cellpitch{0.90}%
    \def\TN@rowgap{0.40em}%
    \def\TN@relationgap{0.56em}%
    \def\TN@physlabeloff{0.56}%
    \def\TN@labelclearance{0.24}%
    \def\TN@annotationclearance{0.28}%
    \def\TN@chainsep{1.18}%
    \def\TN@chainright{3.42}%
    \def\TN@chainellipsisx{2.30}%
    \def\TN@gaugeradius{0.62}%
    \def\TN@gaugediagonal{0.45}%
}

\tikzset{
    tn layout profile/.is choice,
    tn layout profile/normal/.code={\TN@normalprofile},
    tn layout profile/compact/.code={\TN@compactprofile}
}

\newcommand{\TNLayoutProfile}[1]{\tikzset{tn layout profile=#1}}

% Select the sole default at load time.
\TNLayoutProfile{normal}

% ---------- Semantic nodes ----------

\newcommand{\TN@node}[5][]{%
    \TNPlaceObjectInternal{#3}{#4}%
    \tikzset{tn current pic options/.style={#1}}%
    \pic[
        tn pic name={#3},
        tn pic style={#2},
        tn pic contents={#5}
    ] at (#3Placement) {tn atom};%
    \TNEventInternal{atom|picture=\the\TN@pictureid|name=#3|kind=#2}%
}

\newcommand{\TN@tensor}[3][]{%
    \TN@node[#1]{tn tensor node}{#2}{#3}{}%
}

\newcommand{\TN@component}[4][]{%
    \TN@node[#1]{tn component node}{#2}{#3}{$\scriptstyle #4$}%
}

\newcommand{\TN@factor}[4][]{%
    \TN@node[#1]{tn factor node}{#2}{#3}{$\scriptstyle #4$}%
}

\newcommand{\TN@map}[4][]{%
    \TN@node[#1]{tn map node}{#2}{#3}{$\scriptstyle #4$}%
}

\newcommand{\TN@state}[4][]{%
    \TN@node[#1]{tn state node}{#2}{#3}{$\scriptstyle #4$}%
}

\newcommand{\TN@expression}[4][]{%
    \TN@node[#1]{tn expression node}{#2}{#3}{$\scriptstyle #4$}%
}

\newcommand{\TN@insertion}[3][]{%
    \TN@node[#1]{tn insertion node}{#2}{#3}{}%
}

\newcommand{\TN@labeledinsertion}[4][]{%
    \TN@node[#1]{tn insertion node}{#2}{#3}{}%
    \if\relax\detokenize{#4}\relax\else
        \TN@label[anchor=south]{#2.north}{(0,\TN@labelclearance)}{#4}%
    \fi
}

\newcommand{\TN@junction}[3][]{%
    \TN@node[#1]{tn junction node}{#2}{#3}{}%
}

\newcommand{\TN@anonymousjunction}[2][]{%
    \node[tn junction node,#1] at #2 {};%
}

% An external label does not alter the glyph of the object it names.
\newcommand{\TN@label}[4][]{%
    \node[tn label,#1] at ($(#2)+#3$) {$\scriptstyle #4$};%
}

% ---------- Named ports ----------

% Every composable object gives each boundary index a stable name and a type.
% The registry lets contractions reject physical--virtual mismatches.
\newcommand{\TN@registerport}[2]{%
    \TNPortRegisterInternal{#1}{#2}%
    \expandafter\gdef\csname
        TN@porttype@\the\TN@pictureid @#1\endcsname{#2}%
}

\newcommand{\TN@assertporttype}[2]{%
    \@ifundefined{TN@porttype@\the\TN@pictureid @#1}{%
        \PackageError{tnlean-tikz}{Undeclared tensor-network port `#1'}
            {Declare the endpoint with \string\TN@vport, \string\TN@pport,
             or \string\TN@mport before composing it.}%
    }{%
        \edef\TN@actualporttype{\csname
            TN@porttype@\the\TN@pictureid @#1\endcsname}%
        \def\TN@expectedporttype{#2}%
        \ifx\TN@actualporttype\TN@expectedporttype\else
            \PackageError{tnlean-tikz}{Port `#1' has type
                `\TN@actualporttype', not `#2'}
                {Only two virtual ports, two physical ports, or two morphism
                 ports may be joined.}%
        \fi
    }%
}

\newcommand{\TN@assertdifferentports}[2]{%
    \edef\TN@firstendpoint{#1}%
    \edef\TN@secondendpoint{#2}%
    \ifx\TN@firstendpoint\TN@secondendpoint
        \PackageError{tnlean-tikz}{A trace cannot join port `#1' to itself}
            {Choose two distinct typed endpoints for the trace.}%
    \fi
}

% A checked alias retains both the position and the registered type of the
% original port.  It is used when one atom exposes both an object-independent
% role (for example Combined) and a theorem-specific role (for example MPO).
\newcommand{\TN@aliasport}[2]{%
    \@ifundefined{TN@porttype@\the\TN@pictureid @#2}{%
        \PackageError{tnlean-tikz}{Cannot alias undeclared port `#2'}
            {Declare the source port before calling \string\TNPortAlias.}%
    }{%
        \coordinate (#1) at (#2);%
        \edef\TN@aliasedporttype{\csname
            TN@porttype@\the\TN@pictureid @#2\endcsname}%
        \TNAliasRegisterInternal{#1}{#2}{\TN@aliasedporttype}%
        \expandafter\xdef\csname
            TN@porttype@\the\TN@pictureid @#1\endcsname{\TN@aliasedporttype}%
    }%
}

\newcommand{\TN@vport}[3]{%
    \coordinate (#1) at (#2.#3);%
    \TN@registerport{#1}{virtual}%
}

\newcommand{\TN@pport}[3]{%
    \coordinate (#1) at (#2.#3);%
    \TN@registerport{#1}{physical}%
}

\newcommand{\TN@mport}[3]{%
    \coordinate (#1) at (#2.#3);%
    \TN@registerport{#1}{morphism}%
}

% A terminal is an explicitly named external index or morphism endpoint.  It
% is a mathematical boundary object, not an unnamed numerical end of a line.
\newcommand{\TN@vterminal}[2]{%
    \coordinate (#1) at #2;%
    \TN@registerport{#1}{virtual}%
}

\newcommand{\TN@pterminal}[2]{%
    \coordinate (#1) at #2;%
    \TN@registerport{#1}{physical}%
}

\newcommand{\TN@mterminal}[2]{%
    \coordinate (#1) at #2;%
    \TN@registerport{#1}{morphism}%
}

% Annotation endpoints carry no tensor index type.  They name only the support
% of a brace or other explanatory mark.
\newcommand{\TN@annotationterminal}[2]{%
    \coordinate (#1) at #2;%
}

% A port may also lie at a specified fraction of a boundary segment.  This is
% useful for a box carrying several indices on the same side.
\newcommand{\TN@typedbetweenport}[5]{%
    \coordinate (#2) at ($(#3)!#5!(#4)$);%
    \TN@registerport{#2}{#1}%
}

\newcommand{\TN@veastport}[3]{%
    \TN@typedbetweenport{virtual}{#1}{#2.north east}{#2.south east}{#3}%
}

\newcommand{\TN@vwestport}[3]{%
    \TN@typedbetweenport{virtual}{#1}{#2.north west}{#2.south west}{#3}%
}

\newcommand{\TN@vnorthport}[3]{%
    \TN@typedbetweenport{virtual}{#1}{#2.north west}{#2.north east}{#3}%
}

\newcommand{\TN@vsouthport}[3]{%
    \TN@typedbetweenport{virtual}{#1}{#2.south west}{#2.south east}{#3}%
}

\newcommand{\TN@peastport}[3]{%
    \TN@typedbetweenport{physical}{#1}{#2.north east}{#2.south east}{#3}%
}

\newcommand{\TN@pwestport}[3]{%
    \TN@typedbetweenport{physical}{#1}{#2.north west}{#2.south west}{#3}%
}

\newcommand{\TN@pnorthport}[3]{%
    \TN@typedbetweenport{physical}{#1}{#2.north west}{#2.north east}{#3}%
}

\newcommand{\TN@psouthport}[3]{%
    \TN@typedbetweenport{physical}{#1}{#2.south west}{#2.south east}{#3}%
}

% Standard port bundles give every atom the same role-based interface.  The
% names are formed by appending W, E, N, and S to the object name.
\newcommand{\TN@horizontalports}[1]{%
    \TN@vport{#1W}{#1}{west}%
    \TN@vport{#1E}{#1}{east}%
}

\newcommand{\TN@vverticalports}[1]{%
    \TN@vport{#1N}{#1}{north}%
    \TN@vport{#1S}{#1}{south}%
}

\newcommand{\TN@pverticalports}[1]{%
    \TN@pport{#1N}{#1}{north}%
    \TN@pport{#1S}{#1}{south}%
}

\newcommand{\TN@cardinalports}[1]{%
    \TN@horizontalports{#1}%
    \TN@vverticalports{#1}%
}

% The conventional MPO/MPDO signature has horizontal virtual indices and
% vertical physical indices.
\newcommand{\TN@mpoports}[1]{%
    \TN@horizontalports{#1}%
    \TN@pverticalports{#1}%
}

% ---------- Anchored index lines ----------

\newcommand{\TN@connectpoints}[4][]{%
    \draw[#2,#1] (#3) to (#4);%
}

\newcommand{\TN@vconnectpoints}[3][]{%
    \TN@connectpoints[#1]{tn virtual wire}{#2}{#3}%
}

\newcommand{\TN@pconnectpoints}[3][]{%
    \TN@connectpoints[#1]{tn physical wire}{#2}{#3}%
}

% Port-to-port composition is the preferred form when a construction is
% assembled from independently defined atoms.
\newcommand{\TN@vconnectports}[3][]{%
    \TN@assertporttype{#2}{virtual}%
    \TN@assertporttype{#3}{virtual}%
    \TN@vconnectpoints[#1]{#2}{#3}%
    \TNEventInternal{connection|picture=\the\TN@pictureid|type=virtual|from=#2|to=#3|route=straight}%
}

\newcommand{\TN@pconnectports}[3][]{%
    \TN@assertporttype{#2}{physical}%
    \TN@assertporttype{#3}{physical}%
    \TN@pconnectpoints[#1]{#2}{#3}%
    \TNEventInternal{connection|picture=\the\TN@pictureid|type=physical|from=#2|to=#3|route=straight}%
}

% Orthogonal contractions retain typed endpoints.  The order in the command
% name records whether the first segment is horizontal or vertical.
\newcommand{\TN@vconnectportshv}[3][]{%
    \TN@assertporttype{#2}{virtual}%
    \TN@assertporttype{#3}{virtual}%
    \draw[tn virtual wire,rounded corners=3pt,#1] (#2) -| (#3);%
    \TNEventInternal{connection|picture=\the\TN@pictureid|type=virtual|from=#2|to=#3|route=orthogonal-hv}%
}

\newcommand{\TN@vconnectportsvh}[3][]{%
    \TN@assertporttype{#2}{virtual}%
    \TN@assertporttype{#3}{virtual}%
    \draw[tn virtual wire,rounded corners=3pt,#1] (#2) |- (#3);%
    \TNEventInternal{connection|picture=\the\TN@pictureid|type=virtual|from=#2|to=#3|route=orthogonal-vh}%
}

\newcommand{\TN@pconnectportshv}[3][]{%
    \TN@assertporttype{#2}{physical}%
    \TN@assertporttype{#3}{physical}%
    \draw[tn physical wire,rounded corners=3pt,#1] (#2) -| (#3);%
    \TNEventInternal{connection|picture=\the\TN@pictureid|type=physical|from=#2|to=#3|route=orthogonal-hv}%
}

\newcommand{\TN@pconnectportsvh}[3][]{%
    \TN@assertporttype{#2}{physical}%
    \TN@assertporttype{#3}{physical}%
    \draw[tn physical wire,rounded corners=3pt,#1] (#2) |- (#3);%
    \TNEventInternal{connection|picture=\the\TN@pictureid|type=physical|from=#2|to=#3|route=orthogonal-vh}%
}

\newcommand{\TN@mconnectports}[3][]{%
    \TN@assertporttype{#2}{morphism}%
    \TN@assertporttype{#3}{morphism}%
    \TN@connectpoints[#1]{tn morphism arrow}{#2}{#3}%
    \TNEventInternal{connection|picture=\the\TN@pictureid|type=morphism|from=#2|to=#3|route=straight}%
}

\newcommand{\TN@mconnectportshv}[3][]{%
    \TN@assertporttype{#2}{morphism}%
    \TN@assertporttype{#3}{morphism}%
    \draw[tn morphism arrow,rounded corners=3pt,#1] (#2) -| (#3);%
}

\newcommand{\TN@mconnectportsvh}[3][]{%
    \TN@assertporttype{#2}{morphism}%
    \TN@assertporttype{#3}{morphism}%
    \draw[tn morphism arrow,rounded corners=3pt,#1] (#2) |- (#3);%
}

\newcommand{\TN@labeledmorphism}[4]{%
    \TN@mterminal{#1From}{#2}%
    \TN@mterminal{#1To}{#3}%
    \draw[tn morphism arrow] (#1From) --
        node[midway,above,tn label] {$#4$} (#1To);%
}

\newcommand{\TN@labeledcomparison}[4]{%
    \TN@mterminal{#1From}{#2}%
    \TN@mterminal{#1To}{#3}%
    \draw[tn comparison arrow] (#1From) --
        node[midway,above,tn label] {$#4$} (#1To);%
}

% A trace closure is itself composable when its two endpoints are named ports.
% The last two arguments give the horizontal and vertical clearances.
\newcommand{\TN@vtraceportsbelow}[5][]{%
    \TN@assertdifferentports{#2}{#3}%
    \TN@assertporttype{#2}{virtual}%
    \TN@assertporttype{#3}{virtual}%
    \draw[tn virtual trace,#1] (#2) -- ++(-#4,0) -- ++(0,-#5)
        -| ($(#3)+(#4,0)$) -- (#3);%
    \TNEventInternal{connection|picture=\the\TN@pictureid|type=virtual|from=#2|to=#3|route=trace-below}%
}

\newcommand{\TN@vtraceportsabove}[5][]{%
    \TN@assertdifferentports{#2}{#3}%
    \TN@assertporttype{#2}{virtual}%
    \TN@assertporttype{#3}{virtual}%
    \draw[tn virtual trace,#1] (#2) -- ++(-#4,0) -- ++(0,#5)
        -| ($(#3)+(#4,0)$) -- (#3);%
    \TNEventInternal{connection|picture=\the\TN@pictureid|type=virtual|from=#2|to=#3|route=trace-above}%
}

% A vertical pair of ports may be traced around the right of an object.  The
% last two arguments give the horizontal and vertical clearances.
\newcommand{\TN@vtraceportsright}[5][]{%
    \TN@assertdifferentports{#2}{#3}%
    \TN@assertporttype{#2}{virtual}%
    \TN@assertporttype{#3}{virtual}%
    \draw[tn virtual trace,#1] (#2) -- ++(0,#5)
        -| ($(#3)+(#4,-#5)$) -- ++(-#4,0) -- (#3);%
    \TNEventInternal{connection|picture=\the\TN@pictureid|type=virtual|from=#2|to=#3|route=trace-right}%
}

% Physical trace closures have the same compositional interface and geometry
% as virtual trace closures, but their endpoints and strokes retain physical
% index type.  The last two arguments are the horizontal and vertical
% clearances.
\newcommand{\TN@ptraceportsbelow}[5][]{%
    \TN@assertdifferentports{#2}{#3}%
    \TN@assertporttype{#2}{physical}%
    \TN@assertporttype{#3}{physical}%
    \draw[tn physical trace,#1] (#2) -- ++(-#4,0) -- ++(0,-#5)
        -| ($(#3)+(#4,0)$) -- (#3);%
    \TNEventInternal{connection|picture=\the\TN@pictureid|type=physical|from=#2|to=#3|route=trace-below}%
}

\newcommand{\TN@ptraceportsabove}[5][]{%
    \TN@assertdifferentports{#2}{#3}%
    \TN@assertporttype{#2}{physical}%
    \TN@assertporttype{#3}{physical}%
    \draw[tn physical trace,#1] (#2) -- ++(-#4,0) -- ++(0,#5)
        -| ($(#3)+(#4,0)$) -- (#3);%
    \TNEventInternal{connection|picture=\the\TN@pictureid|type=physical|from=#2|to=#3|route=trace-above}%
}

\newcommand{\TN@ptraceportsright}[5][]{%
    \TN@assertdifferentports{#2}{#3}%
    \TN@assertporttype{#2}{physical}%
    \TN@assertporttype{#3}{physical}%
    \draw[tn physical trace,#1] (#2) -- ++(0,#5)
        -| ($(#3)+(#4,-#5)$) -- ++(-#4,0) -- (#3);%
    \TNEventInternal{connection|picture=\the\TN@pictureid|type=physical|from=#2|to=#3|route=trace-right}%
}

\newcommand{\TN@selectedpath}[2][]{%
    \draw[tn selected path,#1] #2;%
}

\newcommand{\TN@factorpath}[2][]{%
    \draw[tn factor region,#1] #2;%
}

\newcommand{\TN@vopenport}[3][]{%
    \TN@assertporttype{#2}{virtual}%
    \coordinate (#2Open) at ($(#2)+(#3)$);%
    \TN@registerport{#2Open}{virtual}%
    \TN@vconnectports[#1]{#2}{#2Open}%
}

\newcommand{\TN@popenport}[3][]{%
    \TN@assertporttype{#2}{physical}%
    \coordinate (#2Open) at ($(#2)+(#3)$);%
    \TN@registerport{#2Open}{physical}%
    \TN@pconnectports[#1]{#2}{#2Open}%
}

% ---------- Semantic regions ----------

\newcommand{\TN@region}[5][]{%
    \begin{scope}[on background layer]
        \node[#2, fit=#4, inner sep=4pt, #1] (#3) {#5};%
    \end{scope}%
}

\newcommand{\TN@groupregion}[4][]{%
    \TN@region[#1]{tn grouping region}{#2}{#3}{#4}%
}

\newcommand{\TN@groupregionabove}[3]{%
    \TN@groupregion{#1}{#2}{}%
    \TN@label[anchor=south]
        {#1.north}{(0,\TN@annotationclearance)}{#3}%
}

\newcommand{\TN@groupregionselectedbelow}[3]{%
    \TN@groupregion{#1}{#2}{}%
    \TN@label[anchor=north,tn region label selected]
        {#1.south}{(0,-\TN@annotationclearance)}{#3}%
}

\makeatother
